\documentclass[11pt,a4paper]{article}
\usepackage[utf8]{inputenc}
\usepackage[italian]{babel}
\usepackage{amsmath}
\usepackage{amssymb}
\usepackage{listings}
\usepackage{xcolor}
\usepackage{hyperref}
\usepackage{geometry}
\usepackage{graphicx}
\usepackage{float}
\usepackage{booktabs}
\usepackage{array}
\usepackage{longtable}

\geometry{margin=2.5cm}

% Configurazione per il codice
\lstset{
    language=JavaScript,
    basicstyle=\ttfamily\small,
    keywordstyle=\color{blue}\bfseries,
    commentstyle=\color{green!60!black},
    stringstyle=\color{red},
    numbers=left,
    numberstyle=\tiny\color{gray},
    stepnumber=1,
    numbersep=5pt,
    backgroundcolor=\color{gray!10},
    frame=single,
    breaklines=true,
    breakatwhitespace=true,
    tabsize=2,
    showspaces=false,
    showstringspaces=false
}

\title{Sistema di Gamification\\
\large Documentazione Tecnica Completa}
\author{ExtraTracker}
\date{\today}

\begin{document}

\maketitle
\tableofcontents
\newpage

\section{Introduzione}

Questo documento fornisce una documentazione completa ed esaustiva del sistema di gamification implementato in ExtraTracker. Il sistema utilizza un meccanismo di punti esperienza (XP), livelli e streak per incentivare l'utilizzo continuativo della piattaforma.

\subsection{Obiettivi del Sistema}

Il sistema di gamification ha i seguenti obiettivi:
\begin{itemize}
    \item \textbf{Engagement}: Motivare gli utenti a utilizzare regolarmente la piattaforma
    \item \textbf{Consistenza}: Promuovere sessioni di studio e lavoro quotidiane attraverso lo streak
    \item \textbf{Qualità}: Premiare la qualità delle risposte e la velocità di apprendimento
    \item \textbf{Progresso}: Fornire feedback visivo sul progresso dell'utente
\end{itemize}

\subsection{Architettura Generale}

Il sistema di gamification è centralizzato nel servizio \texttt{ActivityService}, che gestisce:
\begin{itemize}
    \item Registrazione di tutte le attività utente
    \item Calcolo degli XP guadagnati
    \item Aggiornamento del livello utente
    \item Gestione dello streak (giorni consecutivi di attività)
    \item Statistiche aggregate
\end{itemize}

\section{Componenti Principali}

\subsection{ActivityService}

Il servizio principale per la gestione della gamification si trova in:
\begin{lstlisting}
server/services/activityService.js
\end{lstlisting}

Questo servizio è responsabile di:
\begin{enumerate}
    \item Registrare tutte le attività utente nel database (\texttt{UserActivity})
    \item Calcolare gli XP guadagnati per ogni attività
    \item Aggiornare il livello dell'utente in base agli XP totali
    \item Gestire lo streak in modo atomico (prevenendo race conditions)
    \item Mantenere statistiche aggregate (sessioni di studio, carte riviste, ecc.)
\end{enumerate}

\subsection{Modello Dati Utente}

I dati di gamification sono memorizzati nel modello \texttt{User} nella struttura:

\begin{lstlisting}
gamification: {
    xp: Number,              // XP totali accumulati
    level: Number,           // Livello corrente (calcolato)
    streak: {
        current: Number,     // Streak corrente (giorni consecutivi)
        best: Number,         // Miglior streak mai raggiunto
        lastActivityDate: Date // Ultima data di attività
    },
    stats: {
        totalStudySessions: Number,
        totalFlashcardsReviewed: Number,
        correctAnswers: Number
    }
}
\end{lstlisting}

\section{Sistema di Livelli}

\subsection{Formula di Calcolo del Livello}

Il livello viene calcolato utilizzando una formula quadratica inversa basata sulla radice quadrata:

\begin{equation}
\text{level} = \max\left(1, \left\lfloor \sqrt{\frac{\text{xp}}{100}} \right\rfloor + 1\right)
\end{equation}

Dove:
\begin{itemize}
    \item $\text{xp}$ è l'XP totale accumulato dall'utente
    \item $100$ è la costante \texttt{XP\_LEVEL\_BASE}
    \item $\lfloor \cdot \rfloor$ è la funzione floor (arrotondamento per difetto)
    \item Il risultato è sempre almeno 1 (livello minimo)
\end{itemize}

\subsection{Implementazione}

\begin{lstlisting}
getLevelFromXP(xp) {
    const safeXp = this._toNumber(xp, 0);
    return Math.max(1, Math.floor(Math.sqrt(safeXp / XP_LEVEL_BASE)) + 1);
}
\end{lstlisting}

\subsection{Calcolo XP per il Prossimo Livello}

L'XP necessario per raggiungere il prossimo livello è calcolato come:

\begin{equation}
\text{nextLevelXp} = \max(100, \text{level}^2 \times 100)
\end{equation}

Per il livello 1, il prossimo livello richiede 100 XP. Per livelli superiori, la formula è $\text{level}^2 \times 100$.

\subsection{Esempi di Progressione}

\begin{table}[H]
\centering
\begin{tabular}{@{}ccccc@{}}
\toprule
\textbf{Livello} & \textbf{XP Minimo} & \textbf{XP Massimo} & \textbf{XP per Prossimo} & \textbf{Delta XP} \\
\midrule
1 & 0 & 99 & 100 & 100 \\
2 & 100 & 399 & 400 & 300 \\
3 & 400 & 899 & 900 & 500 \\
4 & 900 & 1599 & 1600 & 700 \\
5 & 1600 & 2499 & 2500 & 900 \\
10 & 8100 & 9999 & 10000 & 1900 \\
20 & 36100 & 39999 & 40000 & 3900 \\
\bottomrule
\end{tabular}
\caption{Progressione dei livelli}
\end{table}

Come si può osservare, la progressione diventa sempre più difficile, richiedendo più XP per ogni livello successivo. Questo crea un sistema di progressione bilanciato che premia la costanza nel tempo.

\section{Sistema di XP (Punti Esperienza)}

\subsection{Attività che Generano XP}

Il sistema riconosce diversi tipi di attività che possono generare XP:

\begin{enumerate}
    \item \textbf{SESSION\_COMPLETE}: Completamento di una sessione di studio
    \item \textbf{GOAL\_COMPLETED}: Completamento di un obiettivo
    \item \textbf{STREAK\_BONUS}: Bonus per streak (attualmente non utilizzato direttamente)
\end{enumerate}

\subsection{Attività che NON Generano XP}

Le seguenti attività vengono registrate ma non generano XP:
\begin{itemize}
    \item \texttt{GOAL\_CREATED}: Creazione di un obiettivo (0 XP)
    \item \texttt{GOAL\_ARCHIVED}: Archiviazione di un obiettivo (0 XP)
    \item \texttt{WORK\_SESSION\_LOGGED}: Registrazione di una sessione di lavoro (0 XP)
    \item \texttt{WORK\_NOTE\_CREATED}: Creazione di una nota di lavoro (0 XP)
    \item \texttt{HABIT\_CHECKIN}: Check-in di un'abitudine (0 XP)
    \item \texttt{HABIT\_SKIPPED}: Saltare un'abitudine (0 XP)
    \item \texttt{PROJECT\_CREATED}: Creazione di un progetto (0 XP)
    \item \texttt{PROJECT\_COMPLETED}: Completamento di un progetto (0 XP)
\end{itemize}

Queste attività vengono comunque registrate per analytics e tracking, ma non contribuiscono direttamente alla progressione di livello.

\section{Calcolo XP per Sessioni di Studio}

\subsection{Formula Completa}

L'XP guadagnato per una sessione di studio è calcolato come:

\begin{equation}
\text{XP}_{\text{sessione}} = \text{Base} + \text{Corrette} + \text{Speed Bonus} + \text{Streak Bonus}
\end{equation}

Dove:
\begin{align}
\text{Base} &= 10 \text{ XP} \\
\text{Corrette} &= \text{correctCount} \times 2 \\
\text{Speed Bonus} &= \max\left(0, \left\lfloor 10 - \frac{\text{timePerCard}}{3} \right\rfloor\right) \\
\text{Streak Bonus} &= \text{metadata.streakBonus} \text{ (attualmente sempre 0)}
\end{align}

E:
\begin{equation}
\text{timePerCard} = \begin{cases}
\frac{\text{timeSpentSeconds}}{\text{totalCards}} & \text{se } \text{totalCards} > 0 \\
0 & \text{altrimenti}
\end{cases}
\end{equation}

\subsection{Implementazione}

\begin{lstlisting}
_calculateSessionXp(metadata) {
    const baseXP = 10;
    const correctCount = this._toNumber(metadata.correctCount, 0);
    const wrongCount = this._toNumber(metadata.wrongCount, 0);
    const timeSpentSeconds = this._toNumber(metadata.timeSpentSeconds, 0);
    const totalCards = this._toNumber(metadata.totalCards, correctCount + wrongCount);

    const timePerCard = totalCards > 0 ? timeSpentSeconds / totalCards : 0;
    const timeBonus = timePerCard > 0
        ? Math.max(0, Math.round(10 - timePerCard / 3))
        : 0;

    const streakBonus = this._toNumber(metadata.streakBonus, 0);

    return baseXP + (correctCount * 2) + timeBonus + streakBonus;
}
\end{lstlisting}

\subsection{Breakdown Dettagliato}

\subsubsection{Base XP}

Ogni sessione di studio completata garantisce \textbf{10 XP base}, indipendentemente dai risultati. Questo premia la semplice azione di completare una sessione.

\subsubsection{XP per Risposte Corrette}

Ogni risposta corretta vale \textbf{2 XP}. Questo incentiva la qualità delle risposte durante lo studio.

\textbf{Esempio}: Se un utente risponde correttamente a 15 carte su 20:
\begin{equation}
\text{XP}_{\text{corrette}} = 15 \times 2 = 30 \text{ XP}
\end{equation}

\subsubsection{Speed Bonus}

Il bonus velocità premia sessioni rapide ma accurate. La formula è:

\begin{equation}
\text{Speed Bonus} = \max\left(0, \left\lfloor 10 - \frac{\text{secondi per carta}}{3} \right\rfloor\right)
\end{equation}

\textbf{Caratteristiche}:
\begin{itemize}
    \item Il bonus massimo è 10 XP (raggiunto con 0 secondi per carta, teoricamente)
    \item Il bonus diminuisce di 1 XP ogni 3 secondi aggiuntivi per carta
    \item Il bonus diventa 0 quando il tempo per carta supera i 30 secondi
    \item Se non ci sono carte, il bonus è 0
\end{itemize}

\textbf{Esempi}:
\begin{table}[H]
\centering
\begin{tabular}{@{}cc@{}}
\toprule
\textbf{Tempo per Carta} & \textbf{Speed Bonus} \\
\midrule
0-2.9s & 10 XP \\
3-5.9s & 9 XP \\
6-8.9s & 8 XP \\
9-11.9s & 7 XP \\
12-14.9s & 6 XP \\
15-17.9s & 5 XP \\
18-20.9s & 4 XP \\
21-23.9s & 3 XP \\
24-26.9s & 2 XP \\
27-29.9s & 1 XP \\
$\geq$ 30s & 0 XP \\
\bottomrule
\end{tabular}
\caption{Speed Bonus in base al tempo per carta}
\end{table}

\subsubsection{Streak Bonus}

Attualmente, lo streak bonus è sempre impostato a 0. Il campo esiste nella struttura dati per future implementazioni, ma non viene calcolato automaticamente.

\subsection{Esempio Completo di Calcolo}

Consideriamo una sessione con:
\begin{itemize}
    \item 20 carte totali
    \item 17 risposte corrette
    \item 3 risposte sbagliate
    \item 120 secondi totali (2 minuti)
    \item Streak bonus: 0
\end{itemize}

\textbf{Calcolo}:
\begin{align}
\text{timePerCard} &= \frac{120}{20} = 6 \text{ secondi} \\
\text{Speed Bonus} &= \left\lfloor 10 - \frac{6}{3} \right\rfloor = \left\lfloor 10 - 2 \right\rfloor = 8 \text{ XP} \\
\text{Corrette} &= 17 \times 2 = 34 \text{ XP} \\
\text{Base} &= 10 \text{ XP} \\
\text{Streak Bonus} &= 0 \text{ XP} \\
\text{XP Totale} &= 10 + 34 + 8 + 0 = \textbf{52 XP}
\end{align}

\section{Calcolo XP per Obiettivi}

\subsection{GOAL\_COMPLETED}

Quando un obiettivo viene completato, l'utente riceve XP. Il valore predefinito è \textbf{25 XP}, ma può essere personalizzato tramite il campo \texttt{metadata.xpEarned}.

\begin{lstlisting}
if (type === 'GOAL_COMPLETED') {
    return this._toNumber(metadata.xpEarned, 25);
}
\end{lstlisting}

\textbf{Nota}: Attualmente, il sistema non calcola dinamicamente l'XP in base alla difficoltà o alla categoria dell'obiettivo. Tutti gli obiettivi completati danno 25 XP per default.

\subsection{STREAK\_BONUS}

Il tipo di attività \texttt{STREAK\_BONUS} può dare XP bonus, con un default di \textbf{5 XP}. Attualmente, questo tipo di attività non viene utilizzato automaticamente nel sistema.

\begin{lstlisting}
if (type === 'STREAK_BONUS') {
    return this._toNumber(metadata.xpEarned, 5);
}
\end{lstlisting}

\section{Sistema di Streak}

\subsection{Definizione}

Lo \textbf{streak} rappresenta il numero di giorni consecutivi in cui l'utente ha completato almeno un'attività che genera XP o viene registrata nel sistema.

\subsection{Logica di Aggiornamento Atomico}

Il sistema utilizza un aggiornamento atomico per prevenire race conditions quando più richieste arrivano simultaneamente. L'implementazione usa operatori MongoDB atomici invece di read-modify-write.

\subsubsection{Algoritmo}

L'algoritmo di aggiornamento dello streak funziona come segue:

\begin{enumerate}
    \item \textbf{Se \texttt{lastActivityDate} non è oggi}:
    \begin{itemize}
        \item \textbf{Se è ieri}: Incrementa lo streak corrente di 1
        \item \textbf{Altrimenti}: Reset dello streak a 1
        \item Aggiorna \texttt{lastActivityDate} a oggi
    \end{itemize}
    \item \textbf{Se \texttt{lastActivityDate} è già oggi}: Nessuna modifica (idempotente)
\end{enumerate}

\subsubsection{Implementazione}

Il metodo \texttt{updateStreakAtomically} esegue due operazioni atomiche sequenziali:

\begin{lstlisting}
async updateStreakAtomically(userId) {
    const today = this._startOfDay(new Date());
    const yesterdayStart = new Date(today.getTime() - MS_PER_DAY);
    
    // OPERAZIONE 1: Incrementa streak se lastActivityDate è ieri
    const incrementResult = await User.updateOne(
        {
            _id: userId,
            $or: [
                { 'gamification.streak.lastActivityDate': { $exists: false } },
                { 'gamification.streak.lastActivityDate': null },
                {
                    $and: [
                        { 'gamification.streak.lastActivityDate': { $gte: yesterdayStart } },
                        { 'gamification.streak.lastActivityDate': { $lt: today } },
                    ],
                },
            ],
        },
        {
            $inc: { 'gamification.streak.current': 1 },
            $set: { 'gamification.streak.lastActivityDate': today },
        }
    );

    // OPERAZIONE 2: Reset streak se lastActivityDate è più vecchio di ieri
    if (incrementResult.matchedCount === 0) {
        const resetResult = await User.updateOne(
            {
                _id: userId,
                'gamification.streak.lastActivityDate': {
                    $exists: true,
                    $ne: null,
                    $lt: yesterdayStart,
                },
            },
            {
                $set: {
                    'gamification.streak.current': 1,
                    'gamification.streak.lastActivityDate': today,
                },
            }
        );
    }
    
    // Aggiorna best streak se necessario
    // ...
}
\end{lstlisting}

\subsection{Best Streak}

Il sistema mantiene anche il \textbf{best streak}, che rappresenta il miglior streak mai raggiunto dall'utente. Questo valore viene aggiornato automaticamente quando lo streak corrente supera il best streak.

\subsection{Attività che Aggiornano lo Streak}

Tutte le attività registrate tramite \texttt{recordActivity} aggiornano automaticamente lo streak, indipendentemente dal fatto che generino XP o meno. Questo include:
\begin{itemize}
    \item Sessioni di studio completate
    \item Obiettivi creati/completati
    \item Sessioni di lavoro registrate
    \item Check-in di abitudini
    \item Creazione di progetti
\end{itemize}

\section{Flusso di Registrazione Attività}

\subsection{Metodo recordActivity}

Il metodo principale per registrare attività è \texttt{recordActivity}, che esegue le seguenti operazioni:

\begin{enumerate}
    \item \textbf{Validazione}: Verifica che il tipo di attività sia valido
    \item \textbf{Arricchimento Metadata}: Aggiunge \texttt{dayOfWeek} e \texttt{timeOfDay}
    \item \textbf{Calcolo XP}: Calcola gli XP guadagnati per questa attività
    \item \textbf{Registrazione}: Crea un record in \texttt{UserActivity}
    \item \textbf{Aggiornamento Utente}:
    \begin{itemize}
        \item Aggiorna XP totale
        \item Ricalcola livello
        \item Aggiorna streak (atomicamente)
        \item Aggiorna statistiche (se applicabile)
    \end{itemize}
    \item \textbf{Ritorno Risultato}: Restituisce se c'è stato un level-up, XP guadagnati e nuovo livello
\end{enumerate}

\subsection{Pattern Observer}

Il sistema utilizza il Pattern Observer per decoupling tra servizi. I servizi emettono eventi invece di chiamare direttamente \texttt{ActivityService}:

\begin{itemize}
    \item \texttt{worklog.created} $\rightarrow$ \texttt{WORK\_SESSION\_LOGGED}
    \item \texttt{goal.completed} $\rightarrow$ \texttt{GOAL\_COMPLETED}
    \item \texttt{session.completed} $\rightarrow$ \texttt{SESSION\_COMPLETE}
\end{itemize}

L'\texttt{ActivitySubscriber} ascolta questi eventi e li converte in chiamate a \texttt{recordActivity}, utilizzando una coda con retry automatico per gestire errori.

\section{Statistiche Aggregate}

\subsection{Statistiche Mantenute}

Il sistema mantiene le seguenti statistiche aggregate per ogni utente:

\begin{itemize}
    \item \textbf{totalStudySessions}: Numero totale di sessioni di studio completate
    \item \textbf{totalFlashcardsReviewed}: Numero totale di carte riviste
    \item \textbf{correctAnswers}: Numero totale di risposte corrette
\end{itemize}

\subsection{Aggiornamento Statistiche}

Le statistiche vengono aggiornate solo per attività di tipo \texttt{SESSION\_COMPLETE}:

\begin{lstlisting}
if (type === 'SESSION_COMPLETE') {
    const correctCount = this._toNumber(enrichedMetadata.correctCount, 0);
    const wrongCount = this._toNumber(enrichedMetadata.wrongCount, 0);
    const totalCards = this._toNumber(enrichedMetadata.totalCards, correctCount + wrongCount);

    nextStats.totalStudySessions += 1;
    nextStats.totalFlashcardsReviewed += totalCards;
    nextStats.correctAnswers += correctCount;
}
\end{lstlisting}

\section{Integrazione Frontend}

\subsection{Componenti UI}

Il frontend include diversi componenti per visualizzare la gamification:

\subsubsection{UserLevelWidget}

Componente che mostra:
\begin{itemize}
    \item Livello corrente
    \item XP totale e XP necessario per il prossimo livello
    \item Barra di progresso verso il prossimo livello
    \item Streak corrente
\end{itemize}

\textbf{Posizione}: \texttt{src/features/gamification/UserLevelWidget.tsx}

\subsubsection{SessionSummaryModal}

Modal che mostra il riepilogo completo di una sessione di studio completata:
\begin{itemize}
    \item Statistiche (corrette, sbagliate, tempo)
    \item Breakdown XP dettagliato (base, corrette, speed bonus, streak bonus)
    \item Progresso verso il prossimo livello
    \item Messaggi celebrativi per performance eccellenti
\end{itemize}

\textbf{Posizione}: \texttt{src/features/gamification/SessionSummaryModal.tsx}

\subsection{API Endpoints}

\subsubsection{GET /api/dashboard/summary}

Restituisce un riepilogo della dashboard inclusa la gamification:

\begin{lstlisting}
gamification: {
    streak: number,        // Streak corrente
    level: number,         // Livello corrente
    xp: number,           // XP totale
    nextLevelXp: number,  // XP necessario per prossimo livello
    progress: number      // Percentuale progresso (0-100)
}
\end{lstlisting}

\section{Considerazioni Tecniche}

\subsection{Atomicità delle Operazioni}

Il sistema utilizza operazioni atomiche MongoDB per:
\begin{itemize}
    \item Aggiornamento dello streak (prevenzione race conditions)
    \item Incremento XP e aggiornamento livello
\end{itemize}

Questo garantisce che anche in caso di richieste simultanee, i dati rimangano consistenti.

\subsection{Gestione Errori}

Il sistema è progettato per essere resiliente:
\begin{itemize}
    \item Se la registrazione di attività fallisce, l'operazione principale (es. completamento sessione) non viene bloccata
    \item Utilizzo di code con retry automatico per attività asincrone
    \item Fallback a chiamata diretta se Redis non è disponibile
\end{itemize}

\subsection{Performance}

Le operazioni di gamification sono ottimizzate per:
\begin{itemize}
    \item Utilizzo di \texttt{validateModifiedOnly} in Mongoose per salvare solo i campi modificati
    \item Query selettive (solo campi necessari)
    \item Aggiornamenti atomici invece di read-modify-write
\end{itemize}

\section{Punti di Integrazione}

\subsection{StudyService}

Il \texttt{StudyService} chiama \texttt{recordActivity} quando una sessione viene completata:

\begin{lstlisting}
const result = await activityService.recordActivity(userId, 'SESSION_COMPLETE', {
    entityId: deckId,
    category: 'study',
    metadata: {
        correctCount,
        wrongCount,
        timeSpentSeconds,
        totalCards,
        deckId,
        mode,
    },
});
\end{lstlisting}

\subsection{GoalService}

Il \texttt{GoalService} registra attività quando:
\begin{itemize}
    \item Un obiettivo viene creato: \texttt{GOAL\_CREATED} (0 XP)
    \item Un obiettivo viene completato: emette evento \texttt{goal.completed} che viene convertito in \texttt{GOAL\_COMPLETED} (25 XP)
\end{itemize}

\subsection{WorkLogService}

Il \texttt{WorkLogService} emette eventi quando:
\begin{itemize}
    \item Un worklog viene creato: \texttt{worklog.created} $\rightarrow$ \texttt{WORK\_SESSION\_LOGGED} (0 XP)
    \item Un worklog viene aggiornato: \texttt{worklog.updated} $\rightarrow$ \texttt{WORK\_SESSION\_UPDATED} (0 XP)
\end{itemize}

\subsection{ProjectService}

Il \texttt{ProjectService} registra attività quando:
\begin{itemize}
    \item Un progetto viene creato: \texttt{PROJECT\_CREATED} (0 XP)
    \item Un progetto viene completato: \texttt{PROJECT\_COMPLETED} (0 XP)
\end{itemize}

\section{Esempi di Utilizzo}

\subsection{Esempio 1: Sessione di Studio}

Un utente completa una sessione di studio con:
\begin{itemize}
    \item 25 carte totali
    \item 22 corrette, 3 sbagliate
    \item 90 secondi totali
\end{itemize}

\textbf{Calcolo}:
\begin{align}
\text{timePerCard} &= \frac{90}{25} = 3.6 \text{ secondi} \\
\text{Speed Bonus} &= \left\lfloor 10 - \frac{3.6}{3} \right\rfloor = \left\lfloor 10 - 1.2 \right\rfloor = 8 \text{ XP} \\
\text{Corrette} &= 22 \times 2 = 44 \text{ XP} \\
\text{Base} &= 10 \text{ XP} \\
\text{XP Totale} &= 10 + 44 + 8 = \textbf{62 XP}
\end{align}

Se l'utente aveva 450 XP (livello 3), dopo questa sessione avrà:
\begin{align}
\text{XP Nuovo} &= 450 + 62 = 512 \text{ XP} \\
\text{Livello Nuovo} &= \left\lfloor \sqrt{\frac{512}{100}} \right\rfloor + 1 = \left\lfloor 2.26 \right\rfloor + 1 = 3
\end{align}

Nessun level-up in questo caso.

\subsection{Esempio 2: Level-Up}

Un utente al livello 2 (400 XP) completa una sessione che gli dà 150 XP:

\begin{align}
\text{XP Nuovo} &= 400 + 150 = 550 \text{ XP} \\
\text{Livello Nuovo} &= \left\lfloor \sqrt{\frac{550}{100}} \right\rfloor + 1 = \left\lfloor 2.35 \right\rfloor + 1 = 3
\end{align}

L'utente passa dal livello 2 al livello 3! Il sistema restituisce \texttt{leveledUp: true}.

\subsection{Esempio 3: Streak}

Un utente con streak di 5 giorni completa un'attività oggi:

\begin{itemize}
    \item Se l'ultima attività era ieri: streak diventa 6
    \item Se l'ultima attività era oggi: streak rimane 6 (idempotente)
    \item Se l'ultima attività era 3 giorni fa: streak viene resettato a 1
\end{itemize}

\section{Future Implementazioni}

\subsection{Possibili Miglioramenti}

Il sistema è progettato per essere estendibile. Possibili miglioramenti futuri:

\begin{enumerate}
    \item \textbf{Streak Bonus Dinamico}: Calcolare automaticamente bonus XP basati sullo streak corrente
    \item \textbf{XP Differenziati per Obiettivi}: Assegnare XP in base alla difficoltà o categoria dell'obiettivo
    \item \textbf{Bonus per Completamento Progetti}: Assegnare XP quando un progetto viene completato
    \item \textbf{Achievement System}: Badge e achievement per milestone specifiche
    \item \textbf{Leaderboard}: Classifiche per livello, XP, streak
    \item \textbf{Multiplier per Orari}: Bonus XP per sessioni in orari specifici
    \item \textbf{Weekend Bonus}: Bonus per attività durante il weekend
\end{enumerate}

\section{Conclusioni}

Il sistema di gamification di ExtraTracker è un sistema completo e ben strutturato che:

\begin{itemize}
    \item Incentiva l'utilizzo regolare della piattaforma
    \item Premia la qualità (risposte corrette) e la velocità (speed bonus)
    \item Utilizza operazioni atomiche per garantire consistenza
    \item È estendibile per future implementazioni
    \item Fornisce feedback visivo chiaro all'utente
\end{itemize}

La progressione quadratica inversa garantisce che il sistema rimanga bilanciato nel lungo termine, mentre lo streak promuove la costanza quotidiana.

\appendix

\section{Codice Completo ActivityService}

\begin{lstlisting}
// File: server/services/activityService.js
// (Vedi file sorgente per codice completo)
\end{lstlisting}

\section{Costanti}

\begin{table}[H]
\centering
\begin{tabular}{@{}ll@{}}
\toprule
\textbf{Costante} & \textbf{Valore} \\
\midrule
\texttt{XP\_LEVEL\_BASE} & 100 \\
\texttt{SESSION\_BASE\_XP} & 10 \\
\texttt{GOAL\_COMPLETED\_XP} & 25 \\
\texttt{STREAK\_BONUS\_XP} & 5 \\
\texttt{XP\_PER\_CORRECT} & 2 \\
\texttt{SPEED\_BONUS\_MAX} & 10 \\
\bottomrule
\end{tabular}
\caption{Costanti del sistema di gamification}
\end{table}

\section{Tipi di Attività}

\begin{longtable}{@{}lp{8cm}l@{}}
\toprule
\textbf{Tipo} & \textbf{Descrizione} & \textbf{XP} \\
\midrule
\texttt{SESSION\_COMPLETE} & Sessione di studio completata & Variabile \\
\texttt{GOAL\_CREATED} & Obiettivo creato & 0 \\
\texttt{GOAL\_COMPLETED} & Obiettivo completato & 25 (default) \\
\texttt{GOAL\_ARCHIVED} & Obiettivo archiviato & 0 \\
\texttt{WORK\_SESSION\_LOGGED} & Sessione di lavoro registrata & 0 \\
\texttt{WORK\_NOTE\_CREATED} & Nota di lavoro creata & 0 \\
\texttt{HABIT\_CHECKIN} & Check-in abitudine & 0 \\
\texttt{HABIT\_SKIPPED} & Abitudine saltata & 0 \\
\texttt{PROJECT\_CREATED} & Progetto creato & 0 \\
\texttt{PROJECT\_COMPLETED} & Progetto completato & 0 \\
\texttt{STREAK\_BONUS} & Bonus streak & 5 (default) \\
\bottomrule
\caption{Tipi di attività e XP associati}
\end{longtable}

\end{document}
